% =============================================================================
% evaluation.tex -- reconciled "Evaluation Overview".
% Owner: Subagent A, reconciled with docs/experiment_design.md + method memo.
%
% MAJOR DIVERGENCES FROM THE PASTED DRAFT (each flagged. Please review):
% [D1] REMOVED the "Proposed Method: Internal Activation Analysis" as a claimed
% contribution. The advisor constraint forbids proposing a novel detector. It
% survives only as an optional exploratory probe in the Discussion, not a result.
% [D2] REPLACED Precision/Recall/F1 as primary metrics with TPR@low-FPR + log-scale
% ROC + AUC (Carlini S&P'22 convention). P/R/F1 retained only as a secondary,
% fixed-threshold view for the benchmark-flagging use case.
% [D3] REMOVED the "fine-tune to force overfitting on PII" positive control. We do not
% train models. Memorization ground truth comes from public Pile membership and
% Pythia's known duplication counts / checkpoints. Enron PII is studied because
% the Enron Emails corpus is itself a Pile subset (i.e., already in training),
% giving controlled PII ground truth without inducing new memorization.
% [CONFIRMED] Enron Emails is present as a component subset of The Pile, so it is
% cited as such.
% [D4] REPLACED k-fold cross-validation (a trained-classifier validation scheme) with
% multi-seed bootstrap confidence intervals + a permutation test, appropriate for
% threshold-free attack metrics.
% =============================================================================

\section{Evaluation Overview}
\label{sec:eval}

\subsection{Threat model and success criteria}
Section~\ref{sec:threat} fixes the parties, the goals G1 through G3, the access tiers,
and the success criteria. This section instantiates that protocol. Each detector is
evaluated at its minimum access tier as graded in Section~\ref{sec:threat}, and success
is defined by the low-false-positive operating point of Section~\ref{sec:concepts} rather
than by average accuracy.

\subsection{Methods under comparison}
We evaluate \emph{existing} detectors only. We propose no new detector. The four
detectors are defined in Section~\ref{sec:concepts}: LOSS/perplexity~\citep{yeom2018privacy},
Min-K\% Prob~\citep{shi2024detecting}, Min-K\%++~\citep{zhang2025minkpp}, and the
zlib-entropy ratio~\citep{carlini2021extracting}. We run Min-K\% and Min-K\%++ at $k=20$.
Two further tests, also defined in Section~\ref{sec:concepts}, operate off the per-item
likelihood axis: corpus-side $n$-gram overlap~\citep{brown2020gpt3}, which we use to
construct ground-truth contamination labels for benchmark items, and the Oren
permutation/exchangeability test~\citep{oren2024proving} at the benchmark level. The
leakage outcome is prefix-continuation extractable memorization under greedy
decoding~\citep{carlini2023quantifying}. On the controlled corpus we additionally measure
regex-detected PII leakage, framed via the PII-leakage games of
\citet{lukas2023pii}. This measurement returned a null at our scale, and we report it as
such rather than as a contribution (Section~\ref{sec:res-extraction}). Related approaches
we deliberately \emph{do not}
evaluate, guided prompting~\citep{golchin2024timetravel}, neighbourhood and shadow-model
reference attacks~\citep{mattern2023neighbourhood,shokri2017membership}, and the
divergence-style extraction of production
models~\citep{nasr2025scalable}, are discussed in Section~\ref{sec:relatedwork}.
An internal-activation probe is reported, if at all, only as exploratory
analysis in the Discussion, not as a contribution.

\subsection{Data}
Table~\ref{tab:datasets} summarizes the corpora and benchmarks used or referenced below.

\paragraph{Models and corpus.} The primary model is the Pythia
suite~\citep{biderman2023pythia}, trained on the public Pile~\citep{gao2020pile}. Its
reconstructible training order, $154$ checkpoints, multiple sizes, and deduplicated
variant provide exact membership ground truth. Our confound-clean split draws members
from a public mirror of the Pile training set and non-members from the Pile validation
set, stratified across the Pile's constituent subsets. We additionally report the
temporally confounded WikiMIA split as a deliberate contrast, and we cite the MIMIR
construction~\citep{duan2024mia} as the methodological precedent rather than using its
released splits. OLMo~\citep{groeneveld2024olmo} on Dolma~\citep{soldaini2024dolma} is a
secondary replication target. The Pile sits within the broader weakly filtered
web-scrape regime, Common Crawl~\citep{commoncrawl} and its filtered derivatives
C4~\citep{raffel2020c4,dodge2021c4} and RedPajama~\citep{weber2024redpajama}, that makes
benchmark contamination structural rather than adversarial.

% =============================================================================
% datasets_table.tex  --  corpora + benchmarks used or referenced in the evaluation.
% Owner: Subagent L. Included by evaluation.tex. Requires \usepackage{booktabs}.
% Every row carries >=1 verified citation (see ../references.bib verification comments).
% Sizes are the figures reported by each artifact's primary source; "approx." where the
% source states an order-of-magnitude or the corpus is continuously growing.
% =============================================================================
\begin{table*}[t]
  \centering
  \caption{Corpora and benchmarks used in or referenced by the evaluation. The Pile is
  our ground-truth training corpus; Common Crawl and C4/Dolma/RedPajama frame the
  weakly filtered web-scrape regime; the lower block lists the contamination benchmarks
  whose items we label by corpus-side overlap.}
  \label{tab:datasets}
  \small
  \begin{tabular}{@{}llp{0.40\linewidth}ll@{}}
    \toprule
    \textbf{Dataset} & \textbf{Type} & \textbf{What it is} & \textbf{Size} & \textbf{Cite} \\
    \midrule
    The Pile      & corpus    & Curated 22-subset English corpus; Pythia's training data and our membership ground truth & 825\,GB & \cite{gao2020pile} \\
    Common Crawl  & corpus    & Open, continually updated repository of raw web-crawl data; the base of most LLM pre-training scrapes & petabyte-scale (growing) & \cite{commoncrawl} \\
    C4            & corpus    & Colossal Clean Crawled Corpus: a filtered Common Crawl snapshot introduced with T5 & $\sim$750\,GB & \cite{raffel2020c4,dodge2021c4} \\
    Dolma         & corpus    & Open pre-training corpus; OLMo's training data (replication target) & 3\,T tokens & \cite{soldaini2024dolma} \\
    RedPajama     & corpus    & Open reproduction of an LLaMA-style pre-training mixture & $\sim$30\,T tokens & \cite{weber2024redpajama} \\
    \midrule
    MMLU          & benchmark & Multiple-choice knowledge/reasoning across 57 subjects & 15{,}908 questions & \cite{hendrycks2021mmlu} \\
    GSM8K         & benchmark & Grade-school multi-step math word problems & 8{,}500 problems & \cite{cobbe2021gsm8k} \\
    HumanEval     & benchmark & Hand-written Python programming problems with unit tests & 164 problems & \cite{chen2021humaneval} \\
    HellaSwag     & benchmark & Adversarially filtered commonsense sentence completion & $\sim$70{,}000 items & \cite{zellers2019hellaswag} \\
    TruthfulQA    & benchmark & Questions probing imitative falsehoods & 817 questions & \cite{lin2022truthfulqa} \\
    BoolQ         & benchmark & Naturally occurring yes/no reading-comprehension questions & 15{,}942 questions & \cite{clark2019boolq} \\
    \bottomrule
  \end{tabular}
\end{table*}


\paragraph{Benchmarks and PII.} Contamination is tested against MMLU, GSM8K, and
HumanEval at the scale reported here. HellaSwag, TruthfulQA, and BoolQ are configured in
the harness and included in Table~\ref{tab:datasets} for reference, and are left for the
GPU replication. For the PII measurement we use the Enron Emails data \emph{as a Pile
subset already present in Pythia's training data}, rather than fine-tuning a model to
memorize PII. The measurement returned no detected leakage at this scale and is reported
as a null (Section~\ref{sec:res-extraction}). No real PII is reproduced in the paper.

\subsection{Metrics}
The rationale for reporting true-positive rate at a low fixed false-positive rate rather
than average accuracy is given in Section~\ref{sec:concepts}. Following that
convention~\citep{carlini2022lira}, the primary metric is \emph{true-positive rate at a
fixed low false-positive rate} (TPR @ $0.1\%$ and $1\%$ FPR) reported with
\emph{log-scale ROC}. AUC-ROC is reported secondarily. For benchmark flagging at a chosen
operating threshold we additionally
report precision/recall/F1 as a secondary, application-facing view. The leakage outcome
is the \emph{extraction rate}~\citep{carlini2023quantifying}. The headline analysis is
the \emph{Spearman correlation between per-item contamination score and per-item
extraction outcome}, with bootstrap confidence intervals and a pre-registered
partial-correlation control that isolates the contribution of raw loss, the quantitative
form of RQ2.

\subsection{Validation and controls}
Robustness is established by bootstrap confidence intervals on TPR@FPR and on the
Spearman correlation, and by a permutation/exchangeability test for benchmark-level
contamination~\citep{oren2024proving}. We run the ablations that preempt the standard
confounds and that are feasible at this scale: deduplicated versus non-deduplicated
Pythia (duplication) and frequency controls (string frequency). The model-size scaling
arm, which asks whether the detector-to-extraction relationship strengthens with scale as
memorization does~\citep{carlini2023quantifying}, is built into the pipeline but not run
here, and is reported as GPU-gated in Section~\ref{sec:limitations}. Differentially private
training~\citep{abadi2016deep,li2022dpllm} is discussed as the mitigation direction
(Section~\ref{sec:dp}), not implemented, since it is a producer-side defense applied at
training time rather than an auditor-side detector.

This section instantiates the protocol of Section~\ref{sec:threat} over concrete methods,
data, and metrics. The empirical results
under this protocol, per-detector TPR at low FPR with log-scale ROC, extraction rates,
and the headline detector-to-extraction correlation with confidence
intervals, are reported in the results section, with every reported number tracing to a
logged harness run.
